\documentclass[11pt]{article}

% --- Packages ---
\usepackage[a4paper,margin=1in]{geometry}
\usepackage{lmodern}
\usepackage[T1]{fontenc}
\usepackage{amsmath,amssymb}
\usepackage{siunitx} % for units and scientific notation
\usepackage{booktabs}
\usepackage{microtype}
\usepackage{caption}
\usepackage{subcaption}
\usepackage{tikz}
\usepackage{pgfplots}
\pgfplotsset{compat=1.18}
\sisetup{detect-all,output-exponent-marker=\mathrm{e},per-mode=symbol}

\title{Time Scales: From Seconds to Planck Time}
\author{}
\date{}

\begin{document}
\maketitle

\section*{Executive Summary}
Time can be measured across an extraordinary range of scales, from everyday
human experiences such as a heartbeat or a blink of an eye (on the order of
seconds) down to the unimaginably brief Planck time (\(\approx 5.39\times10^{-44}\)~s),
which marks the theoretical lower limit where current
physics can meaningfully describe events. Between these extremes lie the SI
subunits of time -- milliseconds, microseconds, nanoseconds, picoseconds,
femtoseconds, attoseconds, zeptoseconds, and yoctoseconds -- each relevant to
different scientific and technological domains. Milliseconds capture human
reaction speeds, microseconds and nanoseconds define electronic and
computational processes, femtoseconds mark chemical bond dynamics, while
attoseconds and beyond probe electron motion and nuclear interactions. Together,
these scales underscore the breadth of phenomena that time measurement
encompasses, bridging human perception, modern engineering, and fundamental
physics.

\section*{Reference Table (Seconds to Yoctoseconds + Planck Time)}
\begin{table}[h!]
  \centering
  \caption{Timescales and examples. Deci/centi are included for completeness but rarely used.}
  \setlength{\tabcolsep}{6pt}
  \begin{tabular}{@{}lll l@{}}
    \toprule
    \textbf{Power of 10} & \textbf{Seconds} & \textbf{SI Name} & \textbf{Example / Context} \\
    \midrule
    \(10^{0}\) & %
    \(\num{1}\) & %
    second (s) & %
    Heartbeat; blink (\(\sim\)0.3--0.4 s) \\ %

    \(10^{-1}\)  & %
    \(\num{1e-1}\) & %
    decisecond (ds) & %
    Typical keystroke interval \\ %

    \(10^{-2}\) & %
    \(\num{1e-2}\) & %
    centisecond (cs) & %
    100~fps frame interval \\ %

    \(10^{-3}\) & %
    \(\num{1e-3}\) & %
    millisecond (ms) & %
    Camera flash; nerve impulse \\ %

    \(10^{-6}\) & %
    \(\num{1e-6}\) & %
    microsecond (\(\mu\)s& %
    MCU instruction; sound travels \(\sim\)\SI{0.34}{mm} \\ %

    \(10^{-9}\)  & %
    \(\num{1e-9}\) & %
    nanosecond (ns) & %
    Light \(\sim\)\SI{30}{cm}; CPU cycle \\

    \(10^{-12}\) & %
    \(\num{1e-12}\) & %
    picosecond (ps) & %
    Molecular vibrations \\

    \(10^{-15}\) & %
    \(\num{1e-15}\) & %
    femtosecond (fs) & %
    Ultrafast laser pulses; bond dynamics \\

    \(10^{-18}\) & %
    \(\num{1e-18}\) & %
    attosecond (as) & %
    Electron motion inside atoms \\

    \(10^{-21}\) & %
    \(\num{1e-21}\) & %
    zeptosecond (zs) & %
    Light crossing H nucleus \(\sim\)2~zs \\

    \(10^{-24}\) & %
    \(\num{1e-24}\) & %
    yoctosecond (ys) & %
    Quark--gluon plasma fluctuations (theoretical) \\

    \addlinespace%

    \(\sim 10^{-44}\) & %
    \(\approx 5.39\times10^{-44}\) & %
    Planck time \(t_p\) & %
    Quantum gravity limit \\

    \bottomrule
  \end{tabular}
\end{table}

\clearpage
\section*{Graphics: Timelines}

% Data we will reuse in both plots:
% Powers and labels (aligned with the table)
% We place coordinates at x = exponent for horizontal, y = exponent for vertical.
% Points: 0, -1, -2, -3, -6, -9, -12, -15, -18, -21, -24, -44 (Planck)
% We'll label with SI names + brief example.

% Helper macro to place labeled marks
\newcommand{\tsmark}[3]{% #1 exponent, #2 name, #3 example
  \addplot+[only marks,mark=*] coordinates {(#1,0)};
  \node[anchor=south,rotate=45] at (axis cs:#1,0) {\small \(\mathbf{10^{#1}}\): #2 --- #3};
}

\subsection*{Horizontal Timeline (log powers on x)}
\begin{figure}[h!]
  \centering
  \begin{tikzpicture}
    \begin{axis}[
      height=7cm,width=\textwidth,
      xmin=-45,xmax=1,
      ymin=-0.2,ymax=0.5,
      axis y line=none,
      axis x line*=bottom,
      xtick={-44,-24,-21,-18,-15,-12,-9,-6,-3,-2,-1,0},
      xticklabels={%
        \(10^{-44}\),%
        \(10^{-24}\),%
        \(10^{-21}\),%
        \(10^{-18}\),%
        \(10^{-15}\),%
        \(10^{-12}\),%
        \(10^{-9}\),%
        \(10^{-6}\),%
        \(10^{-3}\),%
        \(10^{-2}\),%
        \(10^{-1}\),%
        \(10^{0}\)%
      },%
      grid=major,
      grid style={dashed,gray!60},
      tick label style={font=\small},
      label style={font=\small}
    ]
      % Selected labels matching the table
      \tsmark{0}{second}{heartbeat/blink}
      \tsmark{-1}{decisecond}{keystroke}
      \tsmark{-2}{centisecond}{100~fps frame}
      \tsmark{-3}{millisecond}{flash/nerve}
      \tsmark{-6}{microsecond}{MCU instr.}
      \tsmark{-9}{nanosecond}{light \(\sim\)30 cm}
      \tsmark{-12}{picosecond}{molecular vib.}
      \tsmark{-15}{femtosecond}{ultrafast lasers}
      \tsmark{-18}{attosecond}{electron motion}
      \tsmark{-21}{zeptosecond}{H nucleus}
      \tsmark{-24}{yoctosecond}{QGP fluct.}
      \tsmark{-44}{Planck time}{quantum gravity}
    \end{axis}
  \end{tikzpicture}
  \caption{Horizontal timeline of time scales from \(10^{0}\) to \(10^{-44}\) seconds.}
\end{figure}

\clearpage
\subsection*{Vertical Timeline (log powers on y)}
\begin{figure}[h!]
  \centering
  \begin{tikzpicture}
    \begin{axis}[
      height=0.85\textheight,width=11cm,
      ymin=-45,ymax=1,
      xmin=0,xmax=1,
      axis x line=none,
      axis y line*=left,
      ytick={-44,-24,-21,-18,-15,-12,-9,-6,-3,-2,-1,0},
      yticklabels={%
        \(10^{-44}\),%
        \(10^{-24}\),%
        \(10^{-21}\),%
        \(10^{-18}\),%
        \(10^{-15}\),%
        \(10^{-12}\),%
        \(10^{-9}\),%
        \(10^{-6}\),%
        \(10^{-3}\),%
        \(10^{-2}\),%
        \(10^{-1}\),%
        \(10^{0}\)%
      },
      grid=major,
      grid style={dashed,gray!60},
      tick label style={font=\small},
      label style={font=\small}
    ]
      % We'll place points at x=0.2 and labels to the right
      \addplot+[only marks,mark=*] coordinates {%
        (0.2,0)%
        (0.2,-1)%
        (0.2,-2)%
        (0.2,-3)%
        (0.2,-6)%
        (0.2,-9)%
        (0.2,-12)%
        (0.2,-15)%
        (0.2,-18)%
        (0.2,-21)%
        (0.2,-24)%
        (0.2,-44)%
      };%

      \node[anchor=west] at (axis cs:0.25,  0) {\small \textbf{second} --- heartbeat/blink};
      \node[anchor=west] at (axis cs:0.25, -1) {\small decisecond --- keystroke};
      \node[anchor=west] at (axis cs:0.25, -2) {\small centisecond --- 100~fps frame};
      \node[anchor=west] at (axis cs:0.25, -3) {\small \textbf{millisecond} --- flash/nerve};
      \node[anchor=west] at (axis cs:0.25, -6) {\small \textbf{microsecond} --- MCU instr.};
      \node[anchor=west] at (axis cs:0.25, -9) {\small \textbf{nanosecond} --- light \(\sim\)30 cm};
      \node[anchor=west] at (axis cs:0.25,-12) {\small \textbf{picosecond} --- molecular vib.};
      \node[anchor=west] at (axis cs:0.25,-15) {\small \textbf{femtosecond} --- ultrafast lasers};
      \node[anchor=west] at (axis cs:0.25,-18) {\small \textbf{attosecond} --- electron motion};
      \node[anchor=west] at (axis cs:0.25,-21) {\small \textbf{zeptosecond} --- H nucleus};
      \node[anchor=west] at (axis cs:0.25,-24) {\small \textbf{yoctosecond} --- QGP fluct.};
      \node[anchor=west] at (axis cs:0.25,-44) {\small \textbf{Planck time} --- quantum gravity};
    \end{axis}
  \end{tikzpicture}
  \caption{Vertical timeline of time scales from \(10^{0}\) to \(10^{-44}\) seconds.}
\end{figure}

\end{document}

